\section{Method}
\label{sec:method}

\subsection{Overview}

We construct branching agent-memory derivation graphs under controlled
conditions, using a generation harness that enforces a fixed set of
mechanical rules (Section~\ref{sec:harness}), on 30 hand-authored
scenarios (24 poisoned across four \texttt{poison\_form} styles, 6
unpoisoned controls, Section~\ref{sec:corpus}), across three LLMs.
Every derived memory is then labeled against its scenario's
ground-truth semantic target by a validated three-signal automated
pipeline (Section~\ref{sec:labeling}), and blast-radius metrics are
computed against a three-layer ground-truth framework that separates
structural lineage, context exposure, and semantic contamination
(Section~\ref{sec:groundtruth}).

\subsection{Generation harness}
\label{sec:harness}

\textbf{Provenance contract.} Every model call goes through a single
function (\texttt{generate\_trace}, \texttt{src/memoryir/harness.py});
influence edges are recorded at generation time from the scenario's
approved specification, not inferred after the fact---\texttt{STRUCTURAL\_PARENT}
edges mark a write's declared true dependencies, \texttt{CO\_RETRIEVED}
edges mark everything else present in the same shared retrieval
context.

\textbf{Shared retrieval per run.} One retrieval event returns a
\texttt{top\_k} set (true parents first, distractors filling remaining
slots); \emph{all} \texttt{write\_fanout} children produced in that run
derive from the same shared context. This is what makes
\texttt{write\_fanout} actually test over-tainting rather than
relabeling independent single-child runs: naive coarse provenance
would link every retrieved item to every child written in the same
run, producing exactly the false-positive source H1/H2 measure.

\textbf{Focus/background prompt split.} Each derivation call is given
its structural parents as ``primary notes to work from'' and everything
else retrieved as background context explicitly marked off-topic---without
this distinction, an early build showed a ``clean sibling'' write would
spuriously blend in unrelated retrieved content, making it impossible
to construct genuine incidentally-exposed test cases.

\textbf{Depth-continuation rule.} The compromised root's structural
lineage (\texttt{child\_1}) continues across depths (each depth's true
parent is the previous depth's \texttt{child\_1}); sibling writes at
each depth (\texttt{child\_2}, and \texttt{child\_3} when
\texttt{write\_fanout=3}) are fresh, distractor-seeded writes that do
not continue their own lineage, keeping the branching structure
tractable while still exercising co-retrieval fan-out at every depth.
\texttt{child\_3} (present only at \texttt{write\_fanout=3}) has no
scenario-authored structural parent---no scenario spec defines a third
dependent---and is treated as a distractor-only sibling at every depth
by design.

\textbf{Cross-model generation.} Three backends implement an identical
\texttt{derive()} interface so the harness is agnostic to which model
produced a trace: gpt-4o-mini and gpt-4o (Azure OpenAI, same API
shape), and Llama-3.3-70B-Instruct (Azure AI Foundry, a distinct REST
API). The prompt construction is shared verbatim across all three---no
per-model prompt tuning---so cross-model comparisons are not confounded
by differential prompting. All generation calls, and every LLM call in
the labeling pipeline below (\S\ref{sec:labeling}), use
\texttt{temperature=0}. API versions: gpt-4o-mini
\texttt{2025-01-01-preview}, gpt-4o \texttt{2024-12-01-preview} (both
Azure OpenAI); Llama-3.3-70B-Instruct \texttt{2024-05-01-preview}
(Azure AI Foundry). We report API versions rather than Azure
resource/deployment identifiers, which are withheld to avoid
correlating this submission with a specific cloud subscription during
double-blind review.

\subsection{Scenario corpus}
\label{sec:corpus}

\textbf{24 poisoned scenarios}, 6 independently-constructed instances
per \texttt{poison\_form} (\texttt{direct\_instruction},
\texttt{embedded\_fact}, \texttt{authoritative\_framing},
\texttt{multi\_hop\_setup})---the independent construction matters:
\texttt{poison\_form} alone is not a valid sampling unit if each style
is a single hand-written payload merely regenerated across seeds, which
would measure only model stochasticity, not generalization across
independently constructed attack instances of a style.
\texttt{scenario\_id} (not \texttt{poison\_form}, not raw factorial
cells) is therefore the real sampling/inferential unit; bootstrap
confidence intervals resample at this level, stratified by
\texttt{poison\_form}.

\textbf{Oracle ground truth is programmatically generated from a
human-approved specification}---not hand-authored per edge, and not
left to an LLM to decide after the fact. A human approves each
scenario's semantics and declared true-dependency structure (applying
a counterfactual test: if a given parent fact were removed, could the
child still fully express the target semantics from the remaining
facts alone? if yes, that parent is not structurally necessary);
deterministic code then derives every oracle edge and node label from
that approved specification.

\textbf{6 clean/unpoisoned controls}, added specifically as a
false-positive baseline: no adversarial content anywhere, with a
\texttt{semantic\_target} describing a plausible-sounding but
never-injected harmful claim in an analogous domain, verified
non-composable from that scenario's own source facts (no distinctive
identifier overlap, only expected topical-vocabulary overlap).

\texttt{multi\_hop\_setup} \textbf{compositional targets} require
combining two premises that individually establish nothing
harmful---the counterfactual test above genuinely requires two
structural parents only for this style; every other style's targets
are established by a single source fact (a second co-occurring fact is
present but \texttt{CO\_RETRIEVED}, not structurally necessary).

\subsection{Automated labeling}
\label{sec:labeling}

Two independent, deliberately un-conflated label layers:
\begin{enumerate}
\item \textbf{Content-level labels} (\texttt{CARRIES} /
  \texttt{REFERENCES} / \texttt{CLEAN}): does a derived memory's actual
  text assert the injected claim, merely reference it, or show no
  trace? For compositional targets, \texttt{CARRIES} requires asserting
  the full \emph{composed} proposition---juxtaposing component premises
  without composing them is \texttt{REFERENCES}, not \texttt{CARRIES},
  a rule frozen after real generation output surfaced a genuine
  boundary case.
\item \textbf{Structural/oracle labels} (\texttt{STRUCTURAL\_PARENT} /
  \texttt{CO\_RETRIEVED} edges; \texttt{TRUE\_DESCENDANT} /
  \texttt{COEXPOSED} / \texttt{UNRELATED} node reachability): is a
  memory actually causally downstream of the compromised root, per the
  scenario's authored contract? Set at generation time from the
  approved specification, independent of whether the resulting text
  still asserts the claim.
\end{enumerate}

The content-level labeler is LLM-primary: a single LLM judge ---
gpt-4o-mini (Azure OpenAI, API version \texttt{2025-01-01-preview},
\texttt{temperature=0}; the same deployment used for that model's
generation traces, \S\ref{sec:harness}) --- given only the semantic
target and candidate text, never scenario identity, oracle labels, or
marker status, produces the final label, unconditionally. An
independent NLI model
(\texttt{MoritzLaurer/DeBERTa-v3-base-mnli-fever-anli}, a public
HuggingFace checkpoint, run locally) and a second LLM adjudicator ---
a separate call to the \emph{same} gpt-4o-mini deployment, not a
different model --- run alongside it for diagnostic logging
only---an earlier design let
the adjudicator override the primary judge on flagged disagreements,
but validation against a 120-sample human-labeled calibration set
showed this was net-harmful (it fixed 3 wrong labels but broke 5
correct ones, concentrated in compositional-target cases), so the
adjudicator's role was demoted to diagnostic-only. \textbf{Validated
against blind human annotation: Cohen's $\kappa = 0.8699$ and, on an
independent rerun of the identical pipeline, $\kappa = 0.8838$}---both
comfortably clear the pre-registered $\kappa \geq 0.6$ gate; the
rerun's difference from the first (despite zero code changes) reflects
genuine \texttt{temperature=0} API non-determinism, not measurement
error, and both runs agree on the pipeline's one known weakness
(REFERENCES-class recall, concentrated in \texttt{multi\_hop\_setup}'s
compositional boundary). Marker-token detection is deliberately
excluded from the label decision---it is instead a separate H3
measurement (surface-marker survival), and letting it influence the
content label would make the labeler manufacture the correlation H3
exists to discover.

\subsection{Ground truth and metrics}
\label{sec:groundtruth}

\textbf{Three-layer framework}, distinguished because they can
diverge---real generation output produced a case where content merely
co-retrieved into a write's context (not a structural parent) still
surfaced in that write's actual semantic content:
\begin{itemize}
\item \textbf{Structural lineage}---reachability via
  \texttt{STRUCTURAL\_PARENT} edges only (\texttt{TRUE\_DESCENDANT}).
\item \textbf{Context exposure}---reachability via any edge, structural
  or co-retrieved (conservative propagation).
\item \textbf{Semantic contamination}---actual
  \texttt{content\_label == CARRIES}, the ground truth blast-radius
  metrics are anchored to.
\end{itemize}

$B_{true}$ = the set of memories with \texttt{content\_label == CARRIES}
downstream of a compromised root (not structural reachability alone,
per the divergence above). $B_{flagged}$ = whatever a given detection
policy flags, computed \emph{without} access to content labels (a real
detector does not have ground truth):

\begin{equation}
P_{BR} = \frac{|B_{flagged} \cap B_{true}|}{|B_{flagged}|} \qquad
R_{BR} = \frac{|B_{flagged} \cap B_{true}|}{|B_{true}|}
\end{equation}
\begin{equation}
\text{inflation} = \frac{|B_{flagged}|}{|B_{true}|}
\end{equation}

reported as N/A, never coerced to zero, whenever $|B_{true}| = 0$. Four
detection policies are compared: \textbf{structural} (lineage-only),
\textbf{context\_exposure}/\textbf{flat\_transitive} (conservative, any
edge---these are the same policy under two names used by different
hypotheses), \textbf{thresholded} (context-exposure edges pruned by a
per-edge attribution score---operationalized as cosine similarity
between parent/child memory embeddings, since the frozen design named
an \texttt{attribution\_threshold} factor without specifying a concrete
score to threshold), and \textbf{depth\_aware} (conservative
propagation within a bounded window of the compromised root,
structural-only beyond it---an explicit operationalization of the
``one baseline, one proposed method'' containment comparison the
frozen design named without an algorithm). Traversal is implemented as
plain breadth-first search over an in-memory adjacency list (not a live
recursive SQL query), validated against a hand-built fixture with known
answers---including a deliberately cyclic graph confirming
termination---before being run against real data.

\textbf{Containment cost} ($C_{regen}$) counts unique generating runs
that would need to be replayed to regenerate a flagged set (one run
producing several exposed memories counts once, not per-object).

\subsection{Experimental design}

\textbf{Full factorial}: 30 scenarios $\times$ 3 models $\times$ 4
\texttt{top\_k} values (1/3/5/10) $\times$ 3 \texttt{write\_fanout}
values (1/2/3) $\times$ 4 \texttt{derivation\_transform} values
(summarize/paraphrase/refine/continue) $\times$ 5 seeds = 21,600
generated traces. Seeds are a within-scenario noise-reduction measure
(justified by confirmed API non-determinism at \texttt{temperature=0}),
not additional statistical replicates---the true inferential unit
remains \texttt{scenario\_id} (n=24 for the poisoned corpus, 6 per
\texttt{poison\_form} stratum), so seed count does not increase
statistical power. \texttt{depth} (0--5) and
\texttt{attribution\_threshold} are analysis-time factors computed
post-hoc from a trace's persisted edges, not regeneration factors---no
new model calls are needed to sweep them, validated by a
prefix-property check (a depth-$d$ prefix of a depth-5 trace equals a
trace generated with \texttt{max\_depth=d} directly).

\textbf{Reporting hierarchy:} pooled-across-24-scenario estimates are
primary for H1--H4; per-\texttt{poison\_form} panels (n=6 each) are
reported as descriptive/exploratory only, never as independently-powered
comparisons---the effective sample size for a \texttt{poison\_form}-level
claim is 6, not the hundreds of underlying trace rows beneath it.

\textbf{Statistics:} paired bootstrap (10,000 resamples, stratified by
\texttt{poison\_form}) for every reported point estimate and 95\% CI;
Wilcoxon signed-rank for paired policy comparisons at matched
scenarios.
